% !TEX encoding = UTF-8 Unicode
% !TeX program = xelatex
% !TEX TS-program = XeLaTeX

\documentclass[12pt, a4paper]{article}
%\documentclass[12pt, a5paper]{article}
\usepackage{mathtools}
\usepackage{algorithm}
\usepackage[noend]{algpseudocode}
\usepackage{unicode-math}
\usepackage{comment}
\usepackage{wrapfig}
\usepackage{xfrac}
\usepackage{tikz}
\usetikzlibrary{arrows, decorations.markings}

\usepackage[pass]{geometry}
%\geometry{legalpaper, landscape, margin=1in}
\usepackage{fullpage}
\usepackage{setspace}
\usepackage[localise]{xepersian}
\setstretch{1.195}
\settextfont[Scale=1]{XB Niloofar}
\setlatintextfont[Scale=1]{XB Niloofar}%{Junicode} %Junicode
\setdigitfont{XB Niloofar}%{XB Zar}

\setlength{\parindent}{0ex}

\definecolor{answercolor}{HTML}{0080ff}
\newcommand{\vsp}{\vspace{34pt}}
\newcommand{\پاسخ}[2]{#1}
%\thispagestyle{empty}
%\newcommand{\پاسخ}[2]{\color{answercolor}\textbf{پاسخ: }#2\color{black}\vspace{-0.3cm}}

\begin{document}
		\tikzset{
		treenode/.style = {align=center, inner sep=0pt, text centered,
			font=\sffamily},
		arn_n/.style = {treenode, circle, white, font=\sffamily\bfseries, draw=black,
			fill=black, text width=1.5em},% arbre rouge noir, noeud noir
				normal/.style = {circle, white, font=\sffamily\bfseries, draw=black,
			fill=black, text width=1.5em},%
		arn_r/.style = {treenode, circle, red, draw=red, 
			text width=1.5em, very thick},% arbre rouge noir, noeud rouge
		arn_s/.style = {treenode, circle, black, draw=black, 
			text width=1.5em, very thick},% arbre rouge noir, noeud rouge
		arn_x/.style = {treenode, rectangle, draw=black,
			minimum width=0.5em, minimum height=0.5em, very thick}% arbre rouge noir, nil
	}
	\begin{tikzpicture}[remember picture,overlay]	\node[anchor=north,yshift=-1.5cm] at (current page.north) {\rl{به نام خدا}};
	\end{tikzpicture}%

\vspace{-1cm}
{\centering
	\textbf{کوئیز اول درس ساختمان داده‌ها و الگوریتم‌ها}\\
مدرس: دکتر محمد قدسی، مهدی صفرنژاد       - ۲۱ اسفند ۹۹\\
\raggedleft\small
زمان پاسخگویی: ۳۰ دقیقه                   - زمان بارگذاری: ۱۵ دقیقه \\
}
\vspace{-0.5cm}
\پرش‌متوسط\پرخط‌ا\\
در سؤالاتی که توضیح خواسته شده، یک پاسخ کوتاه (حدوداً بین یک تا دو خط) مد نظر است. در سؤالاتی که خواستۀ مسئله توضیحی نیست، جواب نهایی کفایت می‌کند و نوشتن مراحل رسیدن به جواب نیاز نیست.\\

\پاراگراف{سؤال ۱}
می‌خواهیم روی یک رشته $n$ بیتی از ۰ و ۱ الگوریتم زیر را اجرا کنیم. با توجه به این الگوریتم به سؤالات زیر پاسخ دهید.
\begin{latin}
	\begin{algorithm}
		\begin{algorithmic}
			\While{$true$} 
			\State $cnt := 0$
			\For{$ i = 1 \to n-2$}
			\If{$s[i] = 1 \; and \; s[i+1] = 0$}
			\State $s[i] := 1 - s[i]$
			\State $s[i+1] := 1 - s[i + 1]$
			\State $cnt := cnt + 1$
			\EndIf
			\EndFor
			\If{$cnt = 0$}
			\State $break$
			\EndIf
			\EndWhile  \label{roy's loop}
		\end{algorithmic}
	\end{algorithm}
\end{latin}


\begin{itemize}
	\item 
	ثابت کنید این الگوریتم به پایان می‌رسد.
	\item
	پیچیدگی زمانی الگوریتم را در بهترین و بدترین حالت حساب کرده و برای هر کدام یک مثال بزنید.
\end{itemize}

\پاراگراف{سؤال ۲}
پیچیدگی تابع $T(n)$ را در هر قسمت به دست آورید.
\begin{itemize}
	\item $T(n) = 3T(n/3) + \Theta(n(logn)^2)$
	\item $T(n) = 5T(n-1) + \Theta(n^2)$
	\item $T(n) = T(n - \sqrt{n}) + O(n)$
\end{itemize}
\vspace{1pt}
\begin{flushleft}
موفق باشید
\end{flushleft}
\end{document}